\begingroup
\footnotesize
\renewcommand{\arraystretch}{1.16}
\setlength{\tabcolsep}{4pt}
\definecolor{qaQuestionBg}{gray}{0.93}
\definecolor{qaPolicy}{HTML}{6F42C1}
\definecolor{qaGate}{HTML}{D97706}
\definecolor{qaCaveat}{HTML}{B42318}
\newcommand{\qaStatus}[2]{\textcolor{#1}{\textbf{#2}}}
\newcommand{\qaPair}[4]{%
  \rowcolor{qaQuestionBg}%
  \multicolumn{2}{@{}L{0.96\textwidth}@{}}{\textbf{Question.} #3}\\*%
  \textbf{Answer.} & #4\\%
  \addlinespace[0.42em]%
}
\newcommand{\qaSubtable}[1]{%
  \par\Needspace{8\baselineskip}\medskip\noindent\textbf{#1}\par\smallskip%
}

\captionof{table}{Live model-facing questions after resolved setup facts have been
moved into the model text and mapping tables.  Questions are shaded light grey and
answers are unshaded.}
\label{tab:explicit-question-answers}

\qaSubtable{Hydraulic Boundary-Data Provenance}
\begin{longtable}{@{}L{0.18\textwidth} L{0.76\textwidth}@{}}
\toprule
\qaPair{qaGate}{Gate}{\qid{Q006} Which timestamp and time-origin convention produced the active niche curve \(p_E(t)\)?}{The active OGS input is clear: the open niche uses
\path{open_niche_seasonal}, while the square boundary is treated as fixed-pressure
support.  The open query is the source time axis.  The provider should confirm the
source timestamps, time zone or date convention, interpolation/smoothing, mapping
to OGS \(t=0\), and whether the ERT-suggested time shift seen in initial tests also
applies to the boundary curve.}
\bottomrule
\end{longtable}

\qaSubtable{Initial State, Excavation, and Prestress}
\begin{longtable}{@{}L{0.18\textwidth} L{0.76\textwidth}@{}}
\toprule
\qaPair{qaGate}{Gate}{\qid{Q012} Does the initial pressure encode post-excavation drainage near the niche?}{No.  The current input has no heterogeneous initial drainage field from
RH/suction, direct pressure, nearby boreholes, or excavation history.  A
measurement-based initialization needs an accepted source field and date before it
can replace the uniform value.}
\qaPair{qaCaveat}{Caveat}{\qid{Q013} Does the model simulate staged excavation or gradual unloading?}{No.  The computation starts from the already-open 2D geometry, with
\(\uu(\xx,0)=\bm{0}\), no active initial stress tensor, and no active top load.
The elapsed time from excavation to the model start date is not closed by the
checked OGS files.}
\qaPair{qaGate}{Gate}{\qid{Q014} Which prestress magnitude, angle, and excavation reference should be used if prestress is activated?}{Published CD-A modelling sources give stress-magnitude evidence, but no accepted
active prestress tensor for this run.  The water-content modelling paper reports
\(\sigma_h=3\,\mathrm{MPa}\) and \(\sigma_v=7\,\mathrm{MPa}\)
\citep[Sec.~3]{WaterContentEDZ2024}; the GETE modelling paper lists
\(\sig_0=-(5,5,7)\times10^6\,\mathrm{Pa}\)
\citep[modelling Table~2]{Ziefle2022GETE}.  The current XML does not apply either
value, and a usable 2D prestress still needs principal directions or rotation angle,
component order, sign convention, value choice, and an excavation-reference state.}
\bottomrule
\end{longtable}

\qaSubtable{OGS Inputs and Material-Field Questions}
\begin{longtable}{@{}L{0.18\textwidth} L{0.76\textwidth}@{}}
\toprule
\qaPair{qaGate}{Gate}{\qid{Q025} Are mechanical parameters fitted now?}{No.  Orthotropic elasticity, \(\alpha_B=1\), Bishop exponent 1, and
saturation-dependent swelling are active but fixed.  Elastic-parameter fitting
requires a later internal modelling decision, hydro-mechanical residuals, and a
settled prestress/excavation reference state.}
\qaPair{qaGate}{Gate}{\qid{Q027} Is there expert consensus that material fields should evolve in time during the multi-year experiment?}{The current OGS baseline uses time-independent permeability, porosity, retention,
and mechanical material fields.  Expert input is needed on whether clay sealing,
fracture closure/opening, EDZ healing, or other time-dependent material evolution is
physically expected strongly enough to justify adding fitted time-dependent fields
rather than explaining the trends by boundary forcing, saturation state, or spatial
heterogeneity.}
\bottomrule
\end{longtable}

\qaSubtable{Model Time and Output Operators}
\begin{longtable}{@{}L{0.18\textwidth} L{0.76\textwidth}@{}}
\toprule
\qaPair{qaGate}{Gate}{\qid{Q032} How are output times tied to observations or boundary curves?}{The run starts at \(t=0\), identified by comments as 18 September 2019, and ends
at \(1.4\times10^8\,\mathrm{s}\).  Measurements require an explicit
date-to-model-time conversion, tolerance, and interpolation or matching rule.  This
gate is now especially important because initial ERT tests suggest a possible
time-axis shift.}
\bottomrule
\end{longtable}

\qaSubtable{Material and Structural Questions}
\begin{longtable}{@{}L{0.18\textwidth} L{0.76\textwidth}@{}}
\toprule
\qaPair{qaGate}{Gate}{\qid{Q042} Does the same bedding angle control mechanics, permeability, and electrical anisotropy?}{The cited structural evidence supports bedding orientation as important context for
permeability candidates, but it does not prove that Darcy, mechanical, electrical,
and prestress orientations must share one hard angle.  Tensor rotations should
state their bedding-angle, prestress-angle, and coordinate conventions explicitly
rather than hiding the choices inside component values.}
\qaPair{qaGate}{Gate}{\qid{Q044} How should faults, cracks, and EDZ features be classified here?}{They are not active discrete OGS geometry or separate material classes in this
chapter.  They can become structural priors, masks, equivalent fields, or
scenarios only after their source geometry, scale, aperture meaning, and
3D-to-2D projection are defined.}
\bottomrule
\end{longtable}
\endgroup
